\NormalFont\ShortTitle{ΕΣΔΡΑΣ}
{\MT ΕΣΔΡΑΣ

\par }\ChapOne{1}{\PP \VerseOne{1}ΚΑΙ ἐν τῷ πρώτῳ ἔτει Κύρου τοῦ βασιλέως Περσῶν, τοῦ τελεσθῆναι λόγον Κυρίου ἀπὸ στόματος Ἱερεμίου, ἐξήγειρε Κύριος τὸ πνεῦμα Κύρου βασιλέως Περσῶν, καὶ παρήγγειλε φωνὴν ἐν πάσῃ βασιλείᾳ αὐτοῦ, καί γε ἐν γραπτῷ, λέγων,
\par }{\PP \VS{2}Οὕτως εἶπε Κύρος βασιλεὺς Περσῶν, πάσας τὰς βασιλείας τῆς γῆς ἔδωκέ μοι Κύριος ὁ Θεὸς τοῦ οὐρανοῦ, καὶ αὐτὸς ἐπεσκέψατο ἐπʼ ἐμὲ τοῦ οἰκοδομῆσαι οἶκον αὐτῷ ἐν Ἱερουσαλὴμ τῇ ἐν τῇ Ἰουδαίᾳ.
\VS{3}Τίς ἐν ὑμῖν ἀπὸ παντὸς τοῦ λαοῦ αὐτοῦ; καὶ ἔσται ὁ Θεὸς αὐτοῦ μετʼ αὐτοῦ, καὶ ἀναβήσεται εἰς Ἱερουσαλὴμ τὴν ἐν τῇ Ἰουδαίᾳ, καὶ οἰκοδομησάτω τὸν οἶκον Θεοῦ Ἰσραὴλ· αὐτὸς ὁ Θεὸς ὁ ἐν Ἱερουσαλήμ.
\VS{4}Καὶ πᾶς ὁ καταλιπόμενος ἀπὸ πάντων τῶν τόπων οὗ αὐτὸς παροικεῖ ἐκεῖ, καὶ λήψονται αὐτὸν ἄνδρες τοῦ τόπου αὐτοῦ ἐν ἀργυρίῳ, καὶ χρυσίῳ, καὶ ἀποσκευῇ, καὶ κτήνεσι μετὰ τοῦ ἑκουσίου εἰς οἶκον τοῦ Θεοῦ τὸν ἐν Ἱερουσαλήμ.
\par }{\PP \VS{5}Καὶ ἀνέστησαν ἄρχοντες τῶν πατριῶν τῶν Ἰούδα καὶ Βενιαμεὶν, καὶ οἱ ἱερεῖς καὶ οἱ Λευεῖται, πάντων ὧν ἐξήγειρεν ὁ Θεὸς τὸ πνεῦμα αὐτῶν τοῦ ἀναβῆναι οἰκοδομῆσαι τὸν οἶκον Κυρίου τὸν ἐν Ἱερουσαλήμ.
\VS{6}Καὶ πάντες οἱ κυκλόθεν ἐνίσχυσαν ἐν χερσὶν αὐτῶν ἐν σκεύεσιν ἀργυρίου, ἐν χρυσῷ, ἐν ἀποσκευῇ, καὶ ἐν κτήνεσι, καὶ ἐν ξενίοις, πάρεξ τῶν ἑκουσίων.
\par }{\PP \VS{7}Καὶ ὁ βασιλεὺς Κύρος ἐξήνεγκε τὰ σκεύη οἴκου Κυρίου, ἃ ἔλαβε Ναβουχοδονόσορ ἀπὸ Ἱερουσαλὴμ καὶ ἔδωκεν αὐτὰ ἐν οἴκῳ θεοῦ αὐτοῦ·
\VS{8}Καὶ ἐξήνεγκεν αὐτὰ Κύρος ὁ βασιλεὺς Περσῶν ἐπὶ χεῖρα Μιθραδάτου γασβαρηνοῦ, καὶ ἠρίθμησεν αὐτὰ τῷ Σασαβασὰρ τῷ ἄρχοντι τοῦ Ἰούδα.
\VS{9}Καὶ οὗτος ὁ ἀριθμὸς αὐτῶν· ψυκτῆρες χρυσοῖ τριάκοντα, καὶ ψυκτῆρες ἀργυροῖ χίλιοι, παρηλλαγμένα ἐννέα καὶ εἴκοσι, κεφουρῆς χρυσοῖ τριάκοντα, καὶ ἀργυροῖ διπλοῖ τετρακόσια δέκα,
\VS{10}καὶ σκεύη ἕτερα χίλια.
\VS{11}Πάντα τὰ σκεύη τῷ χρυσῷ καὶ τῷ ἀργυρῷ πεντακισχίλια τετρακόσια, τὰ πάντα ἀναβαίνοντα μετὰ Σασαβασὰρ ἀπὸ τῆς ἀποικίας ἐκ Βαβυλῶνος εἰς Ἱερουσαλήμ.

\par }\Chap{2}{\PP \VerseOne{1}Καὶ οὗτοι οἱ υἱοὶ τῆς χώρας οἱ ἀναβαίνοντες ἀπὸ τῆς αἰχμαλωσίας τῆς ἀποικίας, ἧς ἀπῴκισε Ναβουχοδονόσορ βασιλεὺς Βαβυλῶνος εἰς Βαβυλῶνα, καὶ ἐπέστρεψαν εἰς Ἱερουσαλὴμ καὶ Ἰούδα ἀνὴρ εἰς πόλιν αὐτοῦ·
\VS{2}Οἳ ἦλθον μετὰ Ζοροβάβελ, Ἰησοῦς, Νεεμίας, Σαραΐας, Ῥεελίας, Μαρδοχαῖος, Βαλασὰν, Μασφὰρ, Βαγουαὶ, Ῥεοὺμ, Βαανά· ἀνδρῶν ἀριθμὸς λαοῦ Ἰσραήλ.
\par }{\PP \VS{3}Υἱοὶ Φαρὲς, δισχίλιοι ἑκατὸν ἑβδομηκονταδύο.
\par }{\PP \VS{4}Υἱοὶ Σαφατία, τριακόσιοι ἑβδομηκονταδύο.
\par }{\PP \VS{5}Υἱοὶ Ἄρες, ἑπτακόσιοι ἑβδομηκονταπέντε.
\par }{\PP \VS{6}Υἱοὶ Φαὰθ Μωὰβ τοῖς υἱοῖς Ἰησουὲ Ἰωὰβ, δισχίλιοι ὀκτακόσιοι δεκαδύο.
\par }{\PP \VS{7}Υἱοὶ Αἰλὰμ, χίλιοι διακόσιοι πεντηκοντατέσσαρες.
\par }{\PP \VS{8}Υἱοὶ Ζατθουὰ, ἐννακόσιοι τεσσαρακονταπέντε.
\par }{\PP \VS{9}Υἱοὶ Ζακχοὺ, ἑπτακόσιοι ἑξήκοντα.
\par }{\PP \VS{10}Υἱοὶ Βανουὶ, ἑξακόσιοι τεσσαρακονταδύο.
\par }{\PP \VS{11}Υἱοὶ Βαβαῒ, ἑξακόσιοι εἰκοσιτρεῖς.
\par }{\PP \VS{12}Υἱοὶ Ἀσγὰδ, χίλιοι διακόσιοι εἰκοσιδύο.
\par }{\PP \VS{13}Υἱοὶ Ἀδωνικὰμ, ἑξακόσιοι ἑξηκονταέξ.
\par }{\PP \VS{14}Υἱοὶ Βαγουὲ, δισχίλιοι πεντηκονταέξ.
\par }{\PP \VS{15}Υἱοὶ Ἀδδὶν, τετρακόσιοι πεντηκοντατέσσαρες.
\par }{\PP \VS{16}Υἱοὶ Ἀτὴρ τῷ Ἐζεκίᾳ, ἐννενηκονταοκτώ.
\par }{\PP \VS{17}Υἱοὶ Βασσοῦ, τριακόσιοι εἰκοσιτρεῖς.
\par }{\PP \VS{18}Υἱοὶ Ἰωρὰ, ἑκατὸν δεκαδύο.
\par }{\PP \VS{19}Υἱοὶ Ἀσοὺμ, διακόσιοι εἰκοσιτρεῖς.
\par }{\PP \VS{20}Υἱοὶ Γαβὲρ, ἐννενηκονταπέντε.
\par }{\PP \VS{21}Υἱοὶ Βεθλαὲμ, ἑκατὸν εἰκοσιτρεῖς.
\par }{\PP \VS{22}Υἱοὶ Νετωφὰ, πεντηκονταέξ.
\par }{\PP \VS{23}Υἱοὶ Ἀναθὼθ, ἑκατὸν εἰκοσιοκτώ.
\par }{\PP \VS{24}Υἱοὶ Ἀζμὼθ, τεσσαρακοντατρεῖς.
\par }{\PP \VS{25}Υἱοὶ Καριαθιαρὶμ, Χαφιρὰ, καὶ Βηρὼθ, ἑπτακόσιοι τεσσαρακοντατρεῖς.
\par }{\PP \VS{26}Υἱοὶ τῆς Ῥαμὰ καὶ Γαβαὰ, ἑξακόσιοι εἰκοσιεῖς.
\par }{\PP \VS{27}Ἄνδρες Μαχμὰς, ἑκατὸν εἰκοσιδύο.
\par }{\PP \VS{28}Ἄνδρες Βαιθὴλ καὶ Ἀϊὰ, τετρακόσιοι εἰκοσιτρεῖς.
\par }{\PP \VS{29}Υἱοὶ Ναβοῦ, πεντηκονταδύο.
\par }{\PP \VS{30}Υἱοὶ Μαγεβὶς, ἑκατὸν πεντηκονταέξ.
\par }{\PP \VS{31}Υἱοὶ Ἠλαμὰρ, χίλιοι διακόσιοι πεντηκοντατέσσαρες.
\par }{\PP \VS{32}Υἱοὶ Ἠλὰμ, τριακόσιοι εἴκοσι.
\par }{\PP \VS{33}Υἱοὶ Λοδαδὶ καὶ Ὠνὼ, ἑπτακόσιοι εἰκοσιπέντε.
\par }{\PP \VS{34}Υἱοὶ Ἱεριχὼ, τριακόσιοι τεσσαρακονταπέντε.
\par }{\PP \VS{35}Υἱοὶ Σεναὰ, τρισχίλιοι ἑξακόσιοι τριάκοντα.
\par }{\PP \VS{36}Καὶ οἱ ἱερεῖς υἱοὶ Ἰεδουὰ τῷ οἴκῳ Ἰησοῖ, ἐννακόσιοι ἑβδομηκοντατρεῖς.
\VS{37}Υἱοὶ Ἐμμὴρ, χίλιοι πεντηκονταδύο.
\VS{38}Υἱοὶ Φασσοὺρ, χίλιοι διακόσιοι τεσσαρακονταεπτά.
\VS{39}Υἱοὶ Ἠρὲμ, χίλιοι ἑπτά.
\par }{\PP \VS{40}Καὶ οἱ Λευῖται υἱοὶ Ἰησοῦ καὶ Καδμιὴλ τοῖς υἱοῖς Ὠδουΐα, ἑβδομηκοντατέσσαρες.
\par }{\PP \VS{41}Οἱ ᾄδοντες υἱοὶ Ἀσὰφ, ἑκατὸν εἰκοσιοκτώ.
\par }{\PP \VS{42}Υἱοὶ τῶν πυλωρῶν, υἱοὶ Σελλοὺμ, υἱοὶ Ἀτὴρ, υἱοὶ Τελμὼν, υἱοὶ Ἀκοὺβ, υἱοὶ Ἀτιτὰ, υἱοὶ Σωβαῒ, οἱ πάντες ἑκατὸν τριακονταεννέα.
\par }{\PP \VS{43}Οἱ Ναθινὶμ, υἱοὶ Σουθία, υἱοὶ Ἀσουφὰ, υἱοὶ Ταβαὼθ,
\VS{44}υἱοὶ Κάδης, υἱοὶ Σιαὰ, υἱοὶ Φαδὼν,
\VS{45}υἱοὶ Λαβανὼ, υἱοὶ Ἀγαβὰ, υἱοὶ Ἀκοὺβ,
\VS{46}υἱοὶ Ἀγὰβ, υἱοὶ Σελαμὶ, υἱοὶ Ἀνὰν,
\VS{47}υἱοὶ Γεδδὴλ, υἱοὶ Γαὰρ, υἱοὶ Ῥαϊὰ,
\VS{48}υἱοὶ Ῥασὼν, υἱοὶ Νεκωδὰ, υἱοὶ Γαζὲμ,
\VS{49}υἱοὶ Ἀζὼ, υἱοὶ Φασὴ, υἱοὶ Βασὶ,
\VS{50}υἱοὶ Ἀσενὰ, υἱοὶ Μοουνὶμ, υἱοὶ Νεφουσὶμ,
\VS{51}υἱοὶ Βακβοὺκ, υἱοὶ Ἀκουφὰ, υἱοὶ Ἀροὺρ,
\VS{52}υἱοὶ Βασαλὼθ, υἱοὶ Μαουδὰ, υἱοὶ Ἀρσὰ,
\VS{53}υἱοὶ Βαρκὸς, υἱοὶ Σισάρα, υἱοὶ Θεμὰ,
\VS{54}υἱοὶ Νασθιὲ, υἱοὶ Ἀτουφά·
\VS{55}Υἱοὶ δούλων Σαλωμὼν, υἱοὶ Σωταῒ, υἱοὶ Σεφηρὰ, υἱοὶ Φαδουρὰ,
\VS{56}υἱοὶ Ἰεηλὰ, υἱοὶ Δαρκὼν, υἱοὶ Γεδὴλ,
\VS{57}υἱοὶ Σαφατία, υἱοὶ Ἀτὶλ, υἱοὶ Φαχερὰθ, υἱοὶ Ἀσεβωεὶμ, υἱοὶ Ἡμεΐ.
\VS{58}Πάντες οἱ Ναθανὶμ, καὶ υἱοὶ Ἀβδησελμὰ, τριακόσιοι ἐννενηκονταδύο.
\par }{\PP \VS{59}Καὶ οὗτοι οἱ ἀναβάντες ἀπὸ Θελμελὲχ, Θελαρησὰ, Χεροὺβ, Ἡδὰν, Ἐμμὴρ· καὶ οὐκ ἐδυνάσθησαν τοῦ ἀναγγεῖλαι οἶκον πατριᾶς αὐτῶν καὶ σπέρμα αὐτῶν, εἰ ἐξ Ἰσραήλ εἰσιν·
\VS{60}Υἱοὶ Δαλαΐα, υἱοὶ Βουὰ, υἱοὶ Τωβίου, υἱοὶ Νεκωδὰ, ἑξακόσιοι πεντηκονταδύο.
\VS{61}Καὶ ἀπὸ τῶν υἱῶν τῶν ἱερέων υἱοὶ Λαβεία, υἱοὶ Ἀκκοὺς, υἱοὶ Βερζελλαῒ, ὃς ἔλαβεν ἀπὸ τῶν θυγατέρων Βερζελλαῒ τοῦ Γαλααδίτου γυναῖκα, καὶ ἐκλήθη ἐπὶ τῷ ὀνόματι αὐτῶν.
\VS{62}Οὗτοι ἐζήτησαν γραφὴν αὐτῶν οἱ μεθωεσὶμ, καὶ οὐχ εὑρέθησαν, καὶ ἠγχιστεύθησαν ἀπὸ τῆς ἱερατείας.
\VS{63}Καὶ εἶπεν ἀθερσασθὰ αὐτοῖς τοῦ μὴ φαγεῖν ἀπὸ τοῦ ἁγίου τῶν ἁγίων, ἕως ἀναστῇ ἱερεὺς τοῖς φωτίζουσι καὶ τοῖς τελείοις.
\par }{\PP \VS{64}Πᾶσα δὲ ἡ ἐκκλησία ὁμοῦ ὡσεὶ τέσσαρες μυριάδες δισχίλιοι τριακόσιοι ἑξήκοντα,
\VS{65}χωρὶς δούλων αὐτῶν καὶ παιδισκῶν αὐτῶν, οὗτοι ἑπτακισχίλιοι τριακόσιοι τριακονταεπτά· καὶ οὗτοι ᾄδοντες καὶ ᾄδουσαι διακόσιοι.
\VS{66}Ἵπποι αὐτῶν, ἑπτακόσιοι τριακονταέξ· ἡμίονοι αὐτῶν, διακόσιοι τεσσαρακονταπέντε·
\VS{67}Κάμηλοι αὐτῶν, τετρακόσιοι τριακονταπέντε· ὄνοι αὐτῶν, ἑξακισχίλιοι ἑπτακόσιοι εἴκοσι.
\par }{\PP \VS{68}Καὶ ἀπὸ ἀρχόντων πατριῶν ἐν τῷ εἰσελθεῖν αὐτοὺς εἰς οἶκον Κυρίου τὸν ἐν Ἱερουσαλὴμ, ἡκουσιάσαντο εἰς οἶκον τοῦ Θεοῦ, τοῦ στῆσαι αὐτὸν ἐπὶ τὴν ἑτοιμασίαν αὐτοῦ·
\VS{69}ὡς ἡ δύναμις αὐτῶν, ἔδωκαν εἰς θησαυρὸν τοῦ ἔργου χρυσίον καθαρὸν μναὶ ἓξ μυριάδες καὶ χίλιαι, καὶ ἀργυρίου μνὰς πεντακισχιλίας, καὶ κόθωνοι τῶν ἱερέων ἑκατόν.
\par }{\PP \VS{70}Καὶ ἐκάθισαν οἱ ἱερεῖς, καὶ οἱ Λευῖται, καὶ οἱ ἀπὸ τοῦ λαοῦ, καὶ οἱ ᾄδοντες, καὶ οἱ πυλωροὶ, καὶ οἱ Ναθινὶμ ἐν πόλεσιν αὐτῶν, καὶ πᾶς Ἰσραὴλ ἐν πόλεσιν αὐτῶν.

\par }\Chap{3}{\PP \VerseOne{1}Καὶ ἔφθασεν ὁ μὴν ὁ ἕβδομος, καὶ οἱ υἱοὶ Ἰσραὴλ ἐν πόλεσιν αὐτῶν, καὶ συνήχθη ὁ λαὸς ὡς ἀνὴρ εἷς εἰς Ἱερουσαλήμ.
\VS{2}Καὶ ἀνέστη Ἰησοῦς ὁ τοῦ Ἰωσεδὲκ καὶ οἱ ἀδελφοὶ αὐτοῦ ἱερεῖς, καὶ Ζοροβάβελ ὁ τοῦ Σαλαθιὴλ καὶ οἱ ἀδελφοὶ αὐτοῦ, καὶ ᾠκοδόμησαν τὸ θυσιαστήριον Θεοῦ Ἰσραὴλ, τοῦ ἀνενέγκαι ἐπʼ αὐτὸ ὁλοκαυτώσεις, κατὰ τὰ γεγραμμένα ἐν νόμῳ Μωυσῆ ἀνθρώπου τοῦ Θεοῦ.
\par }{\PP \VS{3}Καὶ ἡτοίμασαν τὸ θυσιαστήριον ἐπὶ τὴν ἑτοιμασίαν αὐτοῦ, ὅτι ἐν καταπλήξει ἐπʼ αὐτοὺς ἀπὸ τῶν λαῶν τῶν γαιῶν· καὶ ἀνέβη ἐπʼ αὐτὸ ὁλοκαύτωσις τῷ Κυρίῳ τοπρωῒ καὶ εἰς ἑσπέραν.
\VS{4}Καὶ ἐποίησαν τὴν ἑορτὴν τῶν σκηνῶν κατὰ τὸ γεγραμμένον, καὶ ὁλοκαυτώσεις ἡμέραν ἐν ἡμέρᾳ ἐν ἀριθμῷ ὡς ἡ κρίσις, λόγον ἡμέρας ἐν ἡμέρᾳ αὐτοῦ·
\VS{5}Καὶ μετὰ τοῦτο ὁλοκαυτώσεις ἐνδελεχισμοῦ, καὶ εἰς τὰς νουμηνίας καὶ εἰς πάσας ἑορτὰς τῷ Κυρίῳ τὰς ἡγιασμένας, καὶ παντὶ ἑκουσιαζομένῳ ἑκούσιον τῷ Κυρίῳ.
\VS{6}Ἐν ἡμέρᾳ μιᾷ τοῦ μηνὸς τοῦ ἑβδόμου ἤρξαντο ἀναφέρειν ὁλοκαυτώσεις τῷ Κυρίῳ, καὶ ὁ οἶκος τοῦ Κυρίου οὐκ ἐθεμελιώθη.
\VS{7}Καὶ ἔδωκαν ἀργύριον τοῖς λατόμοις καὶ τοῖς τέκτοσι, καὶ βρώματα καὶ ποτὰ, καὶ ἔλαιον τοῖς Σιδωνίοις καὶ τοῖς Τυρίοις, ἐνέγκαι ξύλα κέδρινα ἀπὸ τοῦ Λιβάνου πρὸς θάλασσαν Ἰόππης, κατʼ ἐπιχώρησιν Κύρου βασιλέως Περσῶν ἐπʼ αὐτούς.
\par }{\PP \VS{8}Καὶ ἐν τῷ ἔτει τῷ δευτέρῳ τοῦ ἐλθεῖν αὐτοὺς εἰς οἶκον τοῦ Θεοῦ ἐν Ἱερουσαλὴμ, ἐν μηνὶ τῷ δευτέρῳ ἤρξατο Ζοροβάβελ ὁ τοῦ Σαλαθιὴλ καὶ Ἰησοῦς ὁ τοῦ Ἰωσεδὲκ, καὶ οἱ κατάλοιποι τῶν ἀδελφῶν αὐτῶν οἱ ἱερεῖς καὶ οἱ Λευῖται, καὶ πάντες οἱ ἐρχόμενοι ἀπὸ τῆς αἰχμαλωσίας εἰς Ἱερουσαλὴμ, καὶ ἔστησαν τοὺς Λευίτας ἀπὸ εἰκοσαετοῦς καὶ ἐπάνω ἐπὶ τοὺς ποιοῦντας τὰ ἔργα ἐν οἴκῳ Κυρίου.
\VS{9}Καὶ ἔστη Ἰησοῦς καὶ οἱ υἱοὶ αὐτοῦ καὶ οἱ ἀδελφοὶ αὐτοῦ, Καδμιὴλ καὶ οἱ υἱοὶ αὐτοῦ υἱοὶ Ἰούδα ἐπὶ τοὺς ποιοῦντας τὰ ἔργα ἐν οἴκῳ τοῦ Θεοῦ· υἱοὶ Ἠναδὰδ, υἱοὶ αὐτῶν καὶ οἱ ἀδελφοὶ αὐτῶν οἱ Λευῖται.
\par }{\PP \VS{10}Καὶ ἐθεμελίωσαν τοῦ οἰκοδομῆσαι τὸν οἶκον Κυρίου· καὶ ἔστησαν οἱ ἱερεῖς ἐστολισμένοι ἐν σάλπιγξι, καὶ οἱ Λευῖται υἱοὶ Ἀσὰφ ἐν κυμβάλοις τοῦ αἰνεῖν τὸν Κύριον ἐπὶ χεῖρας Δαυὶδ βασιλέως Ἰσραήλ.
\VS{11}Καὶ ἀπεκρίθησαν ἐν αἴνῳ καὶ ἀνθομολογήσει τῷ Κυρίῳ, ὅτι ἀγαθὸν, ὅτι εἰς τὸν αἰῶνα τὸ ἔλεος αὐτοῦ ἐπὶ Ἰσραήλ· καὶ πᾶς ὁ λαὸς ἐσήμαινε φωνῇ μεγάλῃ αἰνεῖν τῷ Κυρίῳ ἐπὶ τῇ θεμελιώσει τοῦ οἴκου Κυρίου.
\VS{12}Καὶ πολλοὶ ἀπὸ τῶν ἱερέων καὶ τῶν Λευιτῶν καὶ ἄρχοντες τῶν πατριῶν οἱ πρεσβύτεροι οἳ εἴδοσαν τὸν οἴκον τὸν πρῶτον ἐν θεμελιώσει αὐτοῦ, καὶ τοῦτον τὸν οἶκον ἐν ὀφθαλμοῖς αὐτῶν, ἔκλαιον φωνῇ μεγάλῃ· καὶ ὁ ὄχλος ἐν σημασίᾳ μετʼ εὐφροσύνης τοῦ ὑψῶσαι ᾠδήν.
\VS{13}Καὶ οὐκ ἦν ὁ λαὸς ἐπιγινώσκων φωνὴν σημασίας τῆς εὐφροσύνης ἀπὸ τῆς φωνῆς τοῦ κλαυθμοῦ τοῦ λαοῦ, ὅτι ὁ λαὸς ἐκραύγασε φωνῇ μεγάλῃ, καὶ ἡ φωνὴ ἠκούετο ἕως ἀπὸ μακρόθεν.

\par }\Chap{4}{\PP \VerseOne{1}Καὶ ἤκουσαν οἱ θλίβοντες Ἰούδα καὶ Βενιαμὶν, ὅτι υἱοὶ τῆς ἀποικίας οἰκοδομοῦσιν οἶκον τῷ Κυρίῳ Θεῷ Ἰσραὴλ,
\VS{2}καὶ ἤγγισαν πρὸς Ζοροβάβελ καὶ πρὸς τοὺς ἄρχοντας τῶν πατριῶν, καὶ εἶπον αὐτοῖς, οἰκοδομήσομεν μεθʼ ὑμῶν, ὅτι ὡς ὑμεῖς ἐκζητοῦμεν τῷ Θεῷ ἡμῶν, καὶ αὐτῷ ἡμεῖς θυσιάζομεν ἀπὸ ἡμερῶν Ἀσαραδὰν βασιλέως Ἀσσοὺρ τοῦ ἐνέγκαντος ἡμᾶς ὧδε.
\par }{\PP \VS{3}Καὶ εἶπε πρὸς αὐτοὺς Ζοροβάβελ καὶ Ἰησοῦς καὶ οἱ κατάλοιποι τῶν ἀρχόντων τῶν πατριῶν τοῦ Ἰσραὴλ, οὐχ ἡμῖν καὶ ὑμῖν τοῦ οἰκοδομῆσαι οἶκον τῷ Θεῷ ἡμῶν, ὅτι ἡμεῖς ἐπὶ τοαυτὸ οἰκοδομήσομεν τῷ Κυρίω Θεῷ ἡμῶν, ὡς ἐνετείλατο ἡμῖν Κύρος ὁ βασιλεὺς Περσῶν.
\VS{4}Καὶ ἦν ὁ λαὸς τῆς γῆς ἐκλύων τὰς χεῖρας τοῦ λαοῦ Ἰούδα, καὶ ἐνεπόδιζον αὐτοὺς οἰκοδομεῖν,
\VS{5}καὶ μισθούμενοι ἐπʼ αὐτοὺς βουλευόμενοι τοῦ διασκεδάσαι βουλὴν αὐτῶν πάσας τὰς ἡμέρας Κύρου βασιλέως Περσῶν, καὶ ἕως βασιλείας Δαρείου βασιλείας Περσῶν.
\par }{\PP \VS{6}Καὶ ἐν βασιλείᾳ Ἀσσουήρου, καὶ ἐν ἀρχῇ βασιλείας αὐτοῦ ἔγραψαν ἐπιστολὴν ἐπὶ οἰκοῦντας Ἰούδα καὶ Ἱερουσαλήμ.
\VS{7}Καὶ ἐν ἡμέραις Ἀρθασασθὰ ἔγραψεν ἐν εἰρήνῃ Μιθραδάτῃ Ταβεὴλ καὶ τοῖς λοιποῖς συνδούλοις· πρὸς Ἀρθασασθὰ βασιλέα Περσῶν ἔγραψεν ὁ φορολόγος γραφὴν Συριστὶ καὶ ἡρμηνευμένην.
\VS{8}Ῥεοὺμ βαλτὰμ καὶ Σαμψὰ ὁ γραμματεὺς ἔγραψαν ἐπιστολὴν μίαν κατὰ Ἱερουσαλὴμ τῷ Ἀρθασασθὰ βασιλεῖ.
\VS{9}Τάδε ἔκρινε Ῥεοὺμ βαλτὰμ καὶ Σαμψὰ ὁ γραμματεὺς καὶ οἱ κατάλοιποι σύνδουλοι ἡμῶν, Δειναῖοι, Ἀφαρσαθαχαῖοι, Ταρφαλαῖοι, Ἀφαρσαῖοι, Ἀρχυαῖοι, Βαβυλώνιοι, Σουσαναχαῖοι, Δαυαῖοι,
\VS{10}καὶ οἱ κατάλοιποι ἐθνῶν ὧν ἀπῴκισεν Ἀσσεναφὰρ ὁ μέγας καὶ ὁ τίμιος, καὶ κατῴκισεν αὐτοὺς ἐν πόλεσι τῆς Σομόρων καὶ τὸ κατάλοιπον πέραν τοῦ ποταμοῦ.
\VS{11}Αὕτη ἡ διαταγὴ τῆς ἐπιστολῆς, ἧς ἀπέστειλαν πρὸς αὐτόν· πρὸς Ἀρθασασθὰ βασιλέα παῖδές σου ἄνδρες πέραν τοῦ ποταμοῦ.
\par }{\PP \VS{12}Γνωστὸν ἔστω τῷ βασιλεῖ, ὅτι οἱ Ἰουδαῖοι οἱ ἀναβάντες ἀπὸ σοῦ πρὸς ἡμᾶς ἤλθοσαν εἰς Ἱερουσαλὴμ τὴν πόλιν τὴν ἀποστάτιν καὶ πονηρὰν, ἣν οἰκοδομοῦσι· καὶ τὰ τείχη αὐτῆς κατηρτισμένα εἰσὶ, καὶ θεμελίους αὐτῆς ἀνύψωσαν.
\VS{13}Νῦν οὖν γνωστὸν ἔστω τῷ βασιλεῖ, ὅτι ἐὰν ἡ πόλις ἐκείνη ἀνοικοδομηθῇ, καὶ τὰ τείχη αὐτῆς καταρτισθῶσι, φόροι οὐκ ἔσονταί σοι, οὐδὲ δώσουσι· καὶ τοῦτο βασιλεῖς κακοποιεῖ,
\VS{14}καὶ ἀσχημοσύνην βασιλέως οὐκ ἔξεστιν ἡμῖν ἰδεῖν· διὰ τοῦτο ἐπέμψαμεν καὶ ἐγνωρίσαμεν τῷ βασιλεῖ,
\VS{15}ἵνα ἐπισκέψηται ἐν βίβλῳ ὑπομνηματισμοῦ τῶν πατέρων σου, καὶ εὑρήσεις, καὶ γνώσῃ, ὅτι ἡ πόλις ἐκείνη πόλις ἀποστάτις, καὶ κακοποιοῦσα βασιλεῖς καὶ χώρας, καὶ φυγαδείαι δούλων γίνονται ἐν μέσῳ αὐτῆς ἀπὸ ἡμερῶν αἰῶνος, διὰ ταῦτα ἡ πόλις αὕτη ἠρημώθη.
\VS{16}Γνωρίζομεν οὖν ἡμεῖς τῷ βασιλεῖ, ὅτι ἂν ἡ πόλις ἐκείνη οἰκοδομηθῇ, καὶ τὰ τείχη αὐτῆς καταρτισθῇ, οὐκ ἔστι σοι εἰρήνη.
\par }{\PP \VS{17}Καὶ ἀπέστειλεν ὁ βασιλεὺς πρὸς Ῥεοὺμ βαλτὰμ καὶ Σαμψὰ γραμματέα καὶ τοὺς καταλοίπους συνδούλους αὐτῶν τοὺς οἰκοῦντας ἐν Σαμαρείᾳ καὶ τοὺς καταλοίπους πέραν τοῦ ποταμοῦ, εἰρήνην· καὶ φησὶν,
\VS{18}ὁ φορολόγος ὃν ἀπεστείλατε πρὸς ἡμᾶς, ἐκλήθη ἔμπροσθεν ἐμοῦ·
\VS{19}Καὶ παρʼ ἐμοῦ ἐτέθη γνώμη, καὶ ἐπεσκεψάμεθα καὶ εὕραμεν, ὅτι ἡ πόλις ἐκείνη ἀφʼ ἡμερῶν αἰῶνος ἐπὶ βασιλεῖς ἐπαίρεται, καὶ ἀποστάσεις καὶ φυγαδείαι γίνονται ἐν αὐτῇ·
\VS{20}Καὶ βασιλεῖς ἰσχυροὶ ἐγένοντο ἐν Ἱερουσαλὴμ, καὶ ἐπικρατοῦντες ὅλης τῆς πέραν τοῦ ποταμοῦ, καὶ φόροι πλήρεις καὶ μέρος δίδονται αὐτοῖς.
\VS{21}Καὶ νῦν θέτε γνώμην, καταργῆσαι τοὺς ἄνδρας ἐκείνους, καὶ ἡ πόλις ἐκείνη οὐκ οἰκοδομηθήσεται ἔτι· Ὅπως ἀπὸ τῆς γνώμης
\VS{22}πεφυλαγμένοι ἦτε ἄνεσιν ποιῆσαι περὶ τούτου, μή ποτε πληθυνθῇ ἀφανισμὸς εἰς κακοποίησιν βασιλεῦσι.
\par }{\PP \VS{23}Τότε ὁ φορολόγος τοῦ Ἀρθασασθὰ βασιλέως ἀνέγνω ἐνώπιον Ῥεοὺμ βαλτὰμ καὶ Σαμψὰ γραμματέως καὶ συνδούλων αὐτοῦ· καὶ ἐπορεύθησαν σπουδῇ εἶς Ἱερουσαλὴμ καὶ ἐν Ἰούδα, καὶ κατήργησαν αὐτοὺς ἐν ἵπποις καὶ δυνάμει.
\VS{24}Τότε ἤργησε τὸ ἔργον οἴκου τοῦ Θεοῦ τὸ ἐν Ἱερουσαλήμ· καὶ ἦν ἀργοῦν ἕως δευτέρου ἔτους τῆς βασιλείας Δαρείου βασιλέως Περσῶν.

\par }\Chap{5}{\PP \VerseOne{1}Καὶ προεφήτευσεν Ἀγγαῖος ὁ προφήτης καὶ Ζαχαρίας ὁ τοῦ Ἀδδὼ προφητείαν ἐπὶ τοὺς Ἰουδαίους τοὺς ἐν Ἰούδα καὶ Ἱερουσαλὴμ ἐν ὀνόματι Θεοῦ Ἰσραὴλ ἐπʼ αὐτούς.
\VS{2}Τότε ἀνέστησαν Ζοροβάβελ ὁ τοῦ Σαλαθιὴλ καὶ Ἰησοῦς υἱὸς Ἰωσεδὲκ, καὶ ἤρξαντο οἰκοδομῆσαι τὸν οἶκον τοῦ Θεοῦ τὸν ἐν Ἱερουσαλήμ· καὶ μετʼ αὐτῶν οἱ προφῆται τοῦ Θεοῦ βοηθοῦντες αὐτοῖς.
\par }{\PP \VS{3}Ἐν αὐτῷ τῷ καιρῷ ἦλθεν ἐπʼ αὐτοὺς Θανθαναῒ ἔπαρχος πέραν τοῦ ποταμοῦ, καὶ Σαθαρβουζαναῒ, καὶ οἱ σύνδουλοι αὐτῶν, καὶ τοιάδε εἶπαν αὐτοῖς, τίς ἔθηκεν ὑμῖν γνώμην τοῦ οἰκοδομῆσαι τὸν οἶκον τοῦτον, καὶ τὴν χορηγίαν ταύτην καταρτίσασθαι;
\VS{4}Τότε ταῦτα εἴποσαν αὐτοῖς, τίνα ἐστὶ τὰ ὀνόματα τῶν ἀνδρῶν τῶν οἰκοδομούντων τὴν πόλιν ταύτην;
\VS{5}Καὶ οἱ ὀφθαλμοὶ τοῦ Θεοῦ ἐπὶ τὴν αἰχμαλωσίαν Ἰούδα, καὶ οὐ κατήργησαν αὐτοὺς ἕως γνώμη τῷ Δαρείῳ ἀπηνέχθη· καὶ τότε ἀπεστάλη τῷ φορολόγῳ ὑπὲρ τούτου
\VS{6}διασάφησις ἐπιστολῆς, ἧς ἀπέστειλε Θανθαναῒ, ὁ ἔπαρχος τοῦ πέραν τοῦ ποταμοῦ, καὶ Σαθαρβουζαναῒ καὶ οἱ σύνδουλοι αὐτῶν Ἀφαρσαχαῖοι οἱ ἐν τῷ πέραν τοῦ ποταμοῦ, Δαρείῳ τῷ βασιλεῖ·
\VS{7}Ῥήμασιν ἀπέστειλαν πρὸς αὐτόν· καὶ τάδε γέγραπται ἐν αὐτῷ·
\par }{\PP \VS{8}Δαρείῳ τῷ βασιλεῖ εἰρήνη πᾶσα. Γνωστὸν ἔστω τῷ βασιλεῖ, ὅτι ἐπορεύθημεν εἰς τὴν Ἰουδαίαν χώραν εἰς οἶκον τοῦ Θεοῦ τοῦ μεγάλου, καὶ αὐτὸς οἰκοδομεῖται λίθοις ἐκλεκτοῖς, καὶ ξύλα ἐντίθεται ἐν τοῖς τοίχοις, καὶ τὸ ἔργον ἐκεῖνο ἐπιδέξιον γίνεται, καὶ εὐοδοῦται ἐν ταῖς χερσὶν αὐτῶν.
\VS{9}Τότε ἠρωτήσαμεν τοὺς πρεσβυτέρους ἐκείνους, καὶ οὕτως εἴπαμεν αὐτοῖς, τίς ἔθηκεν ὑμῖν γνώμην τὸν οἶκον τοῦτον οἰκοδομῆσαι, καὶ τὴν χορηγίαν ταύτην καταρτίσασθαι;
\VS{10}Καὶ τὰ ὀνόματα αὐτῶν ἠρωτήσαμεν αὐτοὺς γνωρίσαι σοι, ὥστε γράψαι σοι τὰ ὀνόματα τῶν ἀνδρῶν τῶν ἀρχόντων αὐτῶν.
\VS{11}Καὶ τοιοῦτο τὸ ῥῆμα ἀπεκρίθησαν ἡμῖν, λέγοντες, ἡμεῖς ἐσμὲν δοῦλοι τοῦ Θεοῦ τοῦ οὐρανοῦ καὶ τῆς γῆς, καὶ οἰκοδομοῦμεν τὸν οἶκον ὃς ἦν ᾠκοδομημένος πρὸ τούτου ἔτη πολλὰ, καὶ βασιλεὺς τοῦ Ἰσραὴλ μέγας ᾠκοδόμησεν αὐτὸν, καὶ κατηρτίσατο αὐτὸν αὐτοῖς.
\VS{12}Ἀφότε δὲ παρώργισαν οἱ πατέρες ἡμῶν τὸν Θεὸν τοῦ οὐρανοῦ, ἔδωκεν αὐτοὺς εἰς χεῖρας Ναβουχοδονόσορ βασιλέως Βαβυλῶνος τοῦ Χαλδαίου, καὶ τὸν οἴκον τοῦτον κατέλυσε, καὶ τὸν λαὸν ἀπῴκισεν εἰς Βαβυλῶνα.
\VS{13}Αλλʼ ἐν ἔτει πρώτῳ Κύρου τοῦ βασιλέως, Κύρος ὁ βασιλεὺς ἔθετο γνώμην τὸν οἶκον τοῦ Θεοῦ τοῦτον οἰκοδομηθῆναι·
\VS{14}Καὶ τὰ σκεύη τοῦ οἴκου τοῦ Θεοῦ τὰ χρυσᾶ καὶ τὰ ἀργυρᾶ, ἃ Ναβουχοδονόσορ ἐξήνεγκεν ἀπὸ τοῦ οἴκου τοῦ ἐν Ἱερουσαλὴμ, καὶ ἀπήνεγκεν αὐτὰ εἰς τὸν ναὸν του βασιλέως, ἐξήνεγκεν αὐτὰ Κύρος ὁ βασιλεὺς ἀπὸ τοῦ ναοῦ τοῦ βασιλέως, καὶ ἔδωκε τῷ Σαβανασὰρ τῷ θησαυροφύλακι, τῷ ἐπὶ τοῦ θησαυροῦ,
\VS{15}καὶ εἶπεν αὐτῷ, πάντα τὰ σκεύη λάβε καὶ πορεύου, θὲς αὐτὰ ἐν τῷ οἴκῳ τῷ ἐν Ἱερουσαλὴμ εἰς τὸν τόπον αὐτῶν.
\VS{16}Τότε Σαβανασὰρ ἐκεῖνος ἦλθε καὶ ἔδωκε θεμελίους τοῦ οἴκου τοῦ Θεοῦ ἐν Ἱερουσαλὴμ, καὶ ἀπὸ τότε ἕως τοῦ νῦν ᾠκοδομήθη, καὶ οὐκ ἐτελέσθη.
\VS{17}Καὶ νῦν εἰ ἐπὶ τὸν βασιλέα ἀγαθὸν ἐπισκεπήτω ἐν τῷ οἴκῳ τῆς γάζης τοῦ βασιλέως Βαβυλῶνος, ὅπως γνῷς ὅτι ἀπὸ βασιλέως Κύρου ἐτέθη γνώμη οἰκοδομῆσαι τὸν οἶκον τοῦ Θεοῦ ἐκεῖνον τὸν ἐν Ἱερουσαλήμ· καὶ γνοὺς ὁ βασιλεὺς περὶ τούτου, πεμψάτω πρὸς ἡμᾶς.

\par }\Chap{6}{\PP \VerseOne{1}Τότε Δαρεῖος ὁ βασιλεὺς ἔθηκε γνώμην, καὶ ἐπεσκέψατο ἐν ταῖς βιβλιοθήκαις ὅπου ἡ γάζα κεῖται ἐν Βαβυλῶνι.
\VS{2}Καὶ εὑρέθη ἐν πόλει ἐν τῇ βάρει κεφαλὶς μία, καὶ τοῦτο γεγραμμένον ἐν αὐτῇ ὑπόμνημα.
\par }{\PP \VS{3}Ἐν ἔτει πρώτῳ Κύρου βασιλέως, Κύρος ὁ βασιλεὺς ἔθηκε γνώμην περὶ οἴκου ἱεροῦ Θεοῦ τοῦ ἐν Ἱερουσαλήμ· οἶκος οἰκοδομηθήτω, καὶ τόπος οὗ θυσιάζουσι τὰ θυσιάσματα· καὶ ἔθηκεν ἔπαρμα ὕψος πήχεις ἑξήκοντα, πλάτος αὐτοῦ πήχεων ἑξήκοντα·
\VS{4}Καὶ δόμοι λίθινοι κραταιοὶ τρεῖς, καὶ δόμος ξύλινος εἷς, καὶ ἡ δαπάνη ἐξ οἴκου τοῦ βασιλέως δοθήσεται.
\VS{5}Καὶ τὰ σκεύη οἴκου τοῦ Θεοῦ τὰ ἀργυρᾶ καὶ τὰ χρυσᾶ, ἃ Ναβουχοδονόσορ ἐξήνεγκεν ἀπὸ τοῦ οἴκου τοῦ ἐν Ἱερουσαλὴμ, καὶ ἐκόμισεν εἰς Βαβυλῶνα, καὶ δοθήτω καὶ ἀπελθέτω εἰς τὸν ναὸν τὸν ἐν Ἱερουσαλὴμ ἐπὶ τόπου οὗ ἐτέθη ἐν οἴκῳ τοῦ Θεοῦ.
\par }{\PP \VS{6}Νῦν δώσετε ἔπαρχοι πέραν τοῦ ποταμοῦ Σαθαρβουζαναῒ, καὶ οἱ σύνδουλοι αὐτῶν Ἀφαρσαχαῖοι οἱ ἐν τῷ πέραν τοῦ ποταμοῦ μακρὰν ὄντες ἐκεῖθεν,
\VS{7}νῦν ἄφετε τὸ ἔργον οἴκου τοῦ Θεοῦ· οἱ ἀφηγούμενοι τῶν Ἰουδαίων καὶ οἱ πρεσβύτεροι τῶν Ἰουδαίων οἶκον τοῦ Θεοῦ ἐκεῖνον οἰκοδομείτωσαν ἐπὶ τοῦ τόπου αὐτοῦ.
\VS{8}Καὶ ἀπʼ ἐμοῦ γνώμη ἐτέθη, μή ποτε τὶ ποιήσητε μετὰ τῶν πρεσβυτέρων τῶν Ἰουδαίων τοῦ οἰκοδομηθῆναι οἶκον τοῦ Θεοῦ ἐκεῖνον· καὶ ἀπὸ ὑπαρχόντων βασιλέως τῶν φόρων πέραν τοῦ ποταμοῦ ἐπιμελῶς δαπάνη ἔστω διδομένη τοῖς ἀνδράσιν ἐκείνοις τὸ μὴ καταργηθῆναι.
\VS{9}Καὶ ὃ ἂν ὑστέρημα, καὶ υἱοὺς βοῶν καὶ κριῶν, καὶ ἀμνοὺς εἰς ὁλοκαυτώσεις τῷ Θεῷ τοῦ οὐρανοῦ, πυροὺς, ἅλας, οἶνον, ἔλαιον, κατὰ τὸ ῥῆμα ἱερέων τῶν ἐν Ἱερουσαλὴμ ἔστω διδόμενον αὐτοῖς, ἡμέραν ἐν ἡμέρᾳ, ὃ ἐὰν αἰτήσωσιν,
\VS{10}ἵνα ὦσιν εὐωδίας προσφέροντες τῷ Θεῷ τοῦ οὐρανοῦ, καὶ προσεύχωνται εἰς ζωὴν τοῦ βασιλέως καὶ υἱῶν αὐτοῦ.
\VS{11}Καὶ ἀπʼ ἐμοῦ ἐτέθη γνώμη, ὅτι πᾶς ἄνθρωπος ὃς ἀλλάξει τὸ ῥῆμα τοῦτο, καθαιρεθήσεται ξύλον ἐκ τῆς οἰκίας αὐτοῦ, καὶ ὠρθωμένος πληγήσεται ἐπʼ αὐτοῦ, καὶ ὁ οἶκος αὐτοῦ τὸ κατʼ ἐμὲ ποιηθήσεται.
\VS{12}Καὶ ὁ Θεὸς οὗ κατασκηνοῖ τὸ ὄνομα ἐκεῖ, καταστρέψαι πάντα βασιλέα καὶ λαὸν ὃς ἐκτενεῖ τὴν χεῖρα αὐτοῦ ἀλλάξαι ἢ ἀφανίσαι τὸν οἶκον τοῦ Θεοῦ τὸν ἐν Ἱερουσαλήμ· ἐγὼ Δαρεῖος ἔθηκα γνώμην, ἐπιμελῶς ἔσται.
\par }{\PP \VS{13}Τότε Θανθαναῒ ὁ ἔπαρχος πέραν τοῦ ποταμοῦ, Σαθαρβουζαναῒ, καὶ οἱ σύνδουλοι αὐτοῦ, πρὸς ὃ ἀπέστειλε Δαρεῖος βασιλεὺς, οὕτως ἐποίησαν ἐπιμελῶς.
\VS{14}Καὶ οἱ πρεσβύτεροι τῶν Ἰουδαίων ᾠκοδομοῦσαν καὶ οἱ Λευῖται ἐν προφητείᾳ Ἀγγαίου τοῦ προφήτου, καὶ Ζαχαρίου υἱοῦ Ἀδδώ· καὶ ἀνῳκοδόμησαν καὶ κατηρτίσαντο ἀπὸ γνώμης Θεοῦ Ἰσραὴλ, καὶ ἀπὸ γνώμης Κύρου, καὶ Δαρείου, καὶ Ἀρθασασθὰ βασιλέων Περσῶν.
\par }{\PP \VS{15}Καὶ ἐτέλεσαν τὸν οἶκον τοῦτον ἕως ἡμέρας τρίτης μηνὸς Ἀδὰρ, ὅ ἐστιν ἔτος ἕκτον τῆς βασιλείας Δαρείου τοῦ βασιλέως.
\par }{\PP \VS{16}Καὶ ἐποίησαν οἱ υἱοὶ Ἰσραὴλ, οἱ ἱερεῖς καὶ οἱ Λευῖται, καὶ οἱ κατάλοιποι υἱῶν ἀποικεσίας ἐγκαίνια τοῦ οἴκου τοῦ Θεοῦ ἐν εὐφροσύνῃ.
\VS{17}Καὶ προσήνεγκαν εἰς τὰ ἐγκαίνια τοῦ οἴκου τοῦ Θεοῦ μόσχους ἑκατὸν, κριοὺς διακοσίους, ἀμνοὺς τετρακοσίους, χιμάῤῥους αἰγῶν ὑπὲρ ἁμαρτίας ὑπὲρ παντὸς Ἰσραὴλ δώδεκα εἰς ἀριθμὸν φυλῶν Ἰσραήλ.
\VS{18}Καὶ ἔστησαν τοὺς ἱερεῖς ἐν διαιρέσεσιν αὐτῶν, καὶ τοὺς Λευίτας ἐν μερισμοῖς αὐτῶν, ἐπὶ δουλείας Θεοῦ ἐν Ἱερουσαλὴμ, κατὰ τὴν γραφὴν βίβλου Μωυσῆ.
\par }{\PP \VS{19}Καὶ ἐποίησαν οἱ υἱοὶ τῆς ἀποικεσίας τὸ πάσχα τῇ τεσσαρεσκαιδεκάτῃ τοῦ μηνὸς τοῦ πρώτου.
\VS{20}Ὅτι ἐκαθαρίσθησαν οἱ ἱερεῖς καὶ Λευῖται, ἕως εἷς πάντες καθαροί· καὶ ἔσφαξαν τὸ πάσχα τοῖς πᾶσιν υἱοῖς τῆς ἀποικεσίας καὶ τοῖς ἀδελφοῖς αὐτῶν τοῖς ἱερεῦσι καὶ ἑαυτοῖς.
\VS{21}Καὶ ἔφαγον υἱοὶ Ἰσραὴλ τὸ πάσχα, οἱ ἀπὸ τῆς ἀποικεσίας, καὶ πᾶς ὁ χωριζόμενος τῆς ἀκαθαρσίας ἐθνῶν τῆς γῆς πρὸς αὐτοῦς, τοῦ ἐκζητῆσαι Κύριον Θεὸν Ἰσραήλ.
\VS{22}Καὶ ἐποίησαν τὴν ἑορτὴν τῶν ἀζύμων ἑπτὰ ἡμέρας ἐν εὐφροσύνῃ, ὅτι εὔφρανεν αὐτοὺς Κύριος, καὶ ἐπέστρεψε καρδίαν βασιλέως Ἀσσοὺρ ἐπʼ αὐτοὺς κραταιῶσαι τὰς χεῖρας αὐτῶν ἐν ἔργοις οἴκου τοῦ Θεοῦ Ἰσραήλ.

\par }\Chap{7}{\PP \VerseOne{1}Καὶ μετὰ τὰ ῥήματα ταῦτα ἐν βασιλείᾳ Ἀρθασασθὰ βασιλέως Περσῶν, ἀνέβη Ἔσδρας υἱὸς Σαραίου, υἱοῦ Ἀζαρίου, υἱοῦ Χελκία,
\VS{2}υἱοῦ Σελοὺμ, υἱοῦ Σαδδοὺκ, υἱοῦ Ἀχιτὼβ,
\VS{3}υἱοῦ Σαμαρία, υἱοῦ Ἐσριὰ, υἱοῦ Μαρεὼθ,
\VS{4}υἱοῦ Ζαραΐα, υἱοῦ Ὀζίου, υἱοῦ Βοκκὶ,
\VS{5}υἱοῦ Ἀβισουὲ, υἱοῦ Φινεὲς, υἱοῦ Ἐλεάζαρ, υἱοῦ Ἀαρὼν τοῦ ἱερέως τοῦ πρώτου·
\VS{6}Αὐτὸς Ἔσδρας ἀνέβη ἐκ Βαβυλῶνος, καὶ αὐτὸς γραμματεὺς ταχὺς ἐν νόμῳ Μωυσῆ, ὃν ἔδωκε Κύριος ὁ Θεὸς Ἰσραήλ· καὶ ἔδωκεν αὐτῷ ὁ βασιλεὺς, ὅτι χεὶρ Κυρίου Θεοῦ αὐτοῦ ἐπʼ αὐτὸν ἐν πᾶσιν οἷς ἐζήτει αὐτός.
\VS{7}Καὶ ἀνέβησαν ἀπὸ τῶν υἱῶν Ἰσραὴλ, καὶ ἀπὸ τῶν ἱερέων, καὶ ἀπὸ τῶν Λευιτῶν. καὶ οἱ ᾄδοντες, καὶ οἱ πυλωροὶ, καὶ οἱ Ναθινὶμ, εἰς Ἱερουσαλὴμ ἐν ἔτει ἑβδόμῳ τῷ Ἀρθασασθὰ τῷ βασιλεῖ.
\VS{8}Καὶ ἤλθοσαν εἰς Ἱερουσαλὴμ τῷ μηνὶ τῷ πέμπτῳ, τοῦτο τὸ ἔτος ἕβδομον τῷ βασιλεῖ·
\VS{9}Ὅτι ἐν μιᾷ τοῦ μηνὸς τοῦ πρώτου αὐτὸς ἐθεμελίωσε τὴν ἀνάβασιν τὴν ἀπὸ Βαβυλῶνος· ἐν δὲ τῇ πρώτῃ τοῦ μηνὸς τοῦ πέμπτου ἤλθοσαν εἰς Ἱερουσαλὴμ, ὅτι χεὶρ Θεοῦ αὐτοῦ ἦν ἀγαθὴ ἐπʼ αὐτόν.
\VS{10}ὅτι Ἔσδρας ἔδωκεν ἐν καρδίᾳ αὐτοῦ ζητῆσαι τὸν νόμον, καὶ ποιεῖν καὶ διδάσκειν ἐν Ἰσραὴλ προστάγματα καὶ κρίματα.
\par }{\PP \VS{11}Καὶ αὕτη ἡ διασάφησις τοῦ διατάγματος, οὗ ἔδωκεν Ἀρθασασθὰ τῷ Ἔσδρα τῷ ἱερεῖ τῷ γραμματεῖ βιβλίου λόγων ἐντολῶν Κυρίου καὶ προσταγμάτων αὐτοῦ ἐπὶ τὸν Ἰσραήλ.
\par }{\PP \VS{12}Ἀρθασασθὰ βασιλεὺς βασιλέων Ἔσδρᾳ γραμματεῖ νόμου Κυρίου τοῦ Θεοῦ τοῦ οὐρανοῦ· τετελέσθω λόγος καὶ ἡ ἀπόκρισις.
\VS{13}Ἀπʼ ἐμοῦ ἐτέθη γνώμη, ὅτι πᾶς ὁ ἑχουσιαζόμενος ἐν βασιλείᾳ μου ἀπὸ λαοῦ Ἰσραὴλ καὶ ἱερέων καὶ Λευιτῶν πορευθῆναι εἰς Ἱερουσαλὴμ, μετὰ σοῦ πορευθῆναι.
\VS{14}Ἀπὸ προσώπου τοῦ βασιλέως καὶ τῶν ἑπτὰ συμβούλων ἀπεστάλη ἐπισκέψασθαι ἐπὶ τὴν Ἰουδαίαν καὶ εἰς Ἱερουσαλὴμ νόμῳ Θεοῦ αὐτῶν τῷ ἐν χειρί σου·
\VS{15}Καὶ εἰς οἶκον Κυρίου, ἀργύριον καὶ χρυσίον, ὃ ὁ βασιλεὺς καὶ οἱ σύμβουλοι ἑκουσιάσθησαν τῷ Θεῷ τοῦ Ἰσραὴλ τῷ ἐν Ἱερουσαλὴμ κατασκηνοῦντι.
\VS{16}Καὶ πᾶν ἀργύριον καὶ χρυσίον, ὅ, τι ἐὰν εὕρῃς ἐν πάσῃ χώρᾳ Βαβυλῶνος μετὰ ἑκουσιασμοῦ τοῦ λαοῦ, καὶ ἱερέων τῶν ἑκουσιαζομένων εἰς οἶκον Θεοῦ τὸν ἐν Ἱερουσαλήμ.
\VS{17}Καὶ πάντα προσπορευόμενον τοῦτον ἑτοίμως ἔνταξον ἐν βιβλίῳ τούτῳ, μόσχους, κριοὺς, ἀμνοὺς, καὶ θυσίας αὐτῶν, καὶ σπονδὰς αὐτῶν· καὶ προσοίσεις αὐτὰ ἐπὶ τοῦ θυσιαστηρίου τοῦ οἴκου τοῦ Θεοῦ ὑμῶν τοῦ ἐν Ἱερουσαλήμ.
\VS{18}Καὶ εἴ τι ἐπὶ σὲ καὶ τοὺς ἀδελφούς σου ἀγαθυνθῇ ἐν καταλοίπῳ τοῦ ἀργυρίου καὶ τοῦ χρυσίου ποιῆσαι, ὡς ἀρεστὸν τῷ Θεῷ ὑμῶν ποιήσατε.
\VS{19}Καὶ τὰ σκεύη τὰ διδόμενά σοι εἰς λειτουργίαν οἴκου Θεοῦ, παράδος ἐνώπιον τοῦ Θεοῦ ἐν Ἱερουσαλήμ.
\VS{20}Καὶ κατάλοιπον χρείας οἴκου Θεοῦ σου, ὃ ἂν φανῇ σοι δοῦναι, δώσεις ἀπὸ οἴκων γάζης βασιλέως καὶ ἀπʼ ἐμοῦ.
\par }{\PP \VS{21}Ἐγὼ Ἀρθασασθὰ βασιλεὺς ἔθηκα γνώμην πάσαις ταῖς γάζαις ταῖς ἐν πέρα τοῦ ποταμοῦ, ὅτι πᾶν ὃ ἂν αἰτήσῃ ὑμᾶς Ἔσδρας ὁ ἱερεὺς καὶ γραμματεὺς τοῦ Θεοῦ τοῦ οὐρανοῦ, ἑτοίμως γινέσθω·
\VS{22}ἕως ἀργυρίου ταλάντων ἑκατὸν, καὶ ἕως πυροῦ κόρων ἑκατὸν, καὶ ἕως οἶνου βατῶν ἑκατὸν, καὶ ἕως ἐλαίου βατῶν ἑκατὸν, καὶ ἅλας οὗ οὐκ ἔστι γραφή.
\VS{23}Πᾶν ὅ ἐστιν ἐν γνώμῃ Θεοῦ τοῦ οὐρανοῦ, γινέσθω· προσέχετε μήτις ἐπιχειρήσῃ εἰς τὸν οἶκον Θεοῦ τοῦ οὐρανοῦ, μή ποτε γένηται ὀργὴ ἐπὶ τὴν βασιλείαν τοῦ βασιλέως καὶ τῶν υἱῶν αὐτοῦ.
\VS{24}Καὶ ὑμῖν ἐγνώρισται ἐν πᾶσι τοῖς ἱερεῦσι, καὶ τοῖς Λευίταις, ᾄδουσι, πυλωροῖς, Ναθινὶμ, καὶ λειτουργοῖς οἴκου Θεοῦ τοῦτο, φόρος μὴ ἔστω σοι, οὐκ ἐξουσιάσεις καταδουλοῦσθαι αὐτούς.
\VS{25}Καὶ σὺ Ἔσδρα, ὡς ἡ σοφία τοῦ Θεοῦ ἐν χειρί σου, κατάστησον γραμματεῖς καὶ κριτὰς, ἵνα ὦσι κρίνοντες παντὶ τῷ λαῷ τῷ ἐν πέρα τοῦ ποταμοῦ πᾶσι τοῖς εἰδόσι νόμον τοῦ Θεοῦ σου, καὶ τῷ μὴ εἰδότι γνωριεῖτε.
\par }{\PP \VS{26}Καὶ πὰς ὃς ἂν μὴ ᾖ ποιῶν νόμον τοῦ Θεοῦ καὶ νόμον τοῦ βασιλέως ἑτοίμως, τὸ κρίμα ἔσται γινόμενον ἐξ αὐτοῦ, ἐάν τε εἰς θάνατον, ἐάν τε εἰς παιδείαν, ἐάν τε εἰς ζημίαν τοῦ βίου. ἐάν τε εἰς παράδοσιν.
\par }{\PP \VS{27}Εὐλογητὸς Κύριος ὁ Θεὸς τῶν πατέρων ἡμῶν, ὃς ἔδωκεν ἐν καρδίᾳ τοῦ βασιλέως οὕτως, τοῦ δοξάσαι τὸν οἶκον Κυρίου τὸν ἐν Ἱερουσαλὴμ,
\VS{28}καὶ ἐπʼ ἐμὲ ἔκλινεν ἔλεος ἐν ὀφθαλμοῖς τοῦ βασιλέως καὶ τῶν συμβούλων αὐτοῦ, καὶ πάντων τῶν ἀρχόντων τοῦ βασιλέως, τῶν ἐπῃρμένων· καὶ ἐγὼ ἐκραταιώθην ὡς χεὶρ Θεοῦ ἡ ἀγαθὴ ἐπʼ ἐμέ, καὶ συνῆξα ἀπὸ Ἰσραὴλ ἄρχοντας ἀναβῆναι μετʼ ἐμοῦ.

\par }\Chap{8}{\PP \VerseOne{1}Καὶ οὗτοι οἱ ἄρχοντες πατριῶν αὐτῶν οἱ ὁδηγοὶ ἀναβαίνοντες μετʼ ἐμοῦ ἐν βασιλείᾳ Ἀρθασασθὰ τοῦ βασιλέως Βαβυλῶνος.
\VS{2}Ἀπὸ υἱῶν Φινεὲς, Γηρσών· ἀπὸ υἱῶν Ἰθάμαρ, Δανιήλ· ἀπὸ υἱῶν Δαυὶδ, Ἀττούς.
\VS{3}Ἀπὸ υἱῶν Σαχανία, καὶ ἀπὸ υἱῶν Φόρος, Ζαχαρίας, καὶ μετ αὐτοῦ τὸ συστρεμμα ἑκατὸν καὶ πεντήκοντα.
\VS{4}Ἀπὸ υἱῶν Φαὰθ Μωὰβ, Ἐλιανὰ υἱὸς Σαραΐα, καὶ μετʼ αὐτοῦ διακόσιοι τὰ ἀρσενικά.
\VS{5}Καὶ ἀπὸ υἱῶν Ζαθόης, Σεχενίας υἱὸς Ἀζιὴλ, καὶ μετʼ αὐτοῦ τριακόσια τὰ ἀρσενικά.
\VS{6}Καὶ ἀπὸ τῶν υἱῶν Ἀδὶν, Ὠβὴθ υἱὸς Ἰωνάθαν, καὶ μετʼ αὐτοῦ πεντήκοντα τὰ ἀρσενικά.
\VS{7}Καὶ ἀπὸ υἱῶν Ἠλὰμ, Ἰσαΐας υἱὸς Ἀθελία, καὶ μετʼ αὐτοῦ ἑβδομήκοντα τὰ ἀρσενικά.
\VS{8}Καὶ ἀπὸ υἱῶν Σαφατία, Ζαβαδίας υἱὸς Μιχαὴλ, καὶ μετʼ αὐτοῦ ὀγδοήκοντα τὰ ἀρσενικά.
\VS{9}Καὶ ἀπὸ υἱῶν Ἰωὰβ, Ἀβαδία υἱὸς Ἰεϊὴλ, καὶ μετʼ αὐτοῦ διακόσιοι δεκαοκτὼ τὰ ἀρσενικά.
\VS{10}Καὶ ἀπὸ τῶν υἱῶν Βαανὶ, Σελιμοὺθ υἱὸς Ἰωσεφία, καὶ μετʼ αὐτοῦ ἑκατὸν ἑξήκοντα τὰ ἀρσενικά.
\VS{11}Καὶ ἀπὸ υἱῶν Βαβὶ, Ζαχαρίας υἱὸς Βαβὶ, καὶ μετʼ αὐτοῦ εἰκοσιοκτὼ τὰ ἀρσενικά.
\VS{12}Καὶ ἀπὸ υἱῶν Ἀσγὰδ, Ἰωανὰν υἱὸς Ἀκκατὰν, καὶ μετʼ αὐτοῦ ἑκατὸν δέκα τὰ ἀρσενικά.
\VS{13}Καὶ ἀπὸ υἱῶν Ἀδωνικὰμ ἔσχατοι, καὶ ταῦτα τὰ ὀνόματα αὐτῶν, Ἐλιφαλὰτ, Ἰεὴλ, καὶ Σαμαΐα, καὶ μετʼ αὐτῶν ἑξήκοντα τὰ ἀρσενικά.
\VS{14}Καὶ ἀπὸ υἱῶν Βαγουαῒ, Οὐθαῒ καὶ Ζαβοὺδ, καὶ μετʼ αὐτοῦ ἑβδομήκοντα τὰ ἀρσενικά.
\par }{\PP \VS{15}Καὶ συνῆξα αὐτοὺς πρὸς τὸν ποταμὸν τὸν ἐρχόμενον πρὸς τὸν Εὐὶ, καὶ παρενεβάλομεν ἐκεῖ ἡμέρας τρεῖς· καὶ συνῆκα ἐν τῷ λαῷ καὶ ἐν τοῖς ἱερεῦσι, καὶ ἀπὸ υἱῶν Λευὶ οὐχ εὗρον ἐκεῖ.
\VS{16}Καὶ ἀπέστειλα τῷ Ἐλεάζαρ, τῷ Ἀριὴλ, τῷ Σεμεΐᾳ, καὶ τῷ Ἀλωνὰμ, καὶ τῷ Ἰαρὶβ, καὶ τῷ Ἐλνάθαμ, καὶ τῷ Νάθαν, καὶ τῷ Ζαχαρίᾳ, καὶ τῷ Μεσολλὰμ, καὶ τῷ Ἰωαρὶμ, καὶ τῷ Ἐλνάθαν, συνιέντας.
\VS{17}Καὶ ἐξήνεγκα αὐτοὺς ἐπὶ ἄρχοντας ἐν ἀργυρίῳ τοῦ τόπου, καὶ ἔθηκα ἐν στόματι αὐτῶν λόγους λαλῆσαι πρὸς τοὺς ἀδελφοὺς αὐτῶν τῶν Ἀθινεὶμ ἐν ἀργυρίῳ τοῦ τόπου, τοῦ ἐνέγκαι ἡμῖν ᾄδοντας εἰς οἶκον Θεοῦ ἡμῶν.
\VS{18}Καὶ ἤλθοσαν ἡμῖν ὡς χεὶρ Θεοῦ ἡμῶν ἀγαθὴ ἐφʼ ἡμᾶς, ἀνὴρ σαχὼν ἀπὸ υἱῶν Μοολὶ, υἱοῦ Λευὶ, υἱοῦ Ἰσραήλ· καὶ ἀρχὴν ἦλθον οἱ υἱοὶ αὐτοῦ καὶ ἀδελφοὶ αὐτοῦ δεκαοκτώ·
\VS{19}Καὶ τὸν Ἀσεβία, καὶ τὸν Ἰσαΐα ἀπὸ τῶν υἱῶν Μεραρὶ, ἀδελφοὶ αὐτοῦ καὶ υἱοὶ αὐτοῦ εἴκοσι.
\VS{20}Καὶ ἀπὸ τῶν Ναθινὶμ, ὧν ἔδωκε Δαυὶδ καὶ οἱ ἄρχοντες εἰς δουλείαν τῶν Λευιτῶν, Ναθινὶμ διακόσιοι εἴκοσι, πάντες συνήχθησαν ἐν ὀνόμασι.
\par }{\PP \VS{21}Καὶ ἐκάλεσα ἐκεῖ νηστείαν ἐπὶ τὸν ποταμὸν Ἀουὲ, τοῦ ταπεινωθῆναι ἐνώπιον τοῦ Θεοῦ ἡμῶν, ζητῆσαι παρʼ αὐτοῦ ὁδὸν εὐθείαν ἡμῖν καὶ τοῖς τέκνοις ἡμῶν καὶ πάσῃ τῇ κτήσει ἡμῶν·
\VS{22}Ὅτι ᾐσχύνθην αἰτήσασθαι παρὰ τοῦ βασιλέως δύναμιν καὶ ἱππεῖς σῶσαι ἡμᾶς ἀπὸ ἐχθροῦ ἐν τῇ ὁδῷ, ὅτι εἴπαμεν τῷ βασιλεῖ, λέγοντες, χεὶρ τοῦ Θεου ἡμῶν ἐπὶ πάντας τοὺς ζητοῦντας αὐτὸν εἰς ἀγαθόν· καὶ κράτος αὐτοῦ, καὶ θυμὸς αὐτοῦ, ἐπὶ πάντας τοὺς ἐγκαταλείποντας αὐτόν.
\VS{23}Καὶ ἐνηστεύσαμεν, καὶ ἐζητήσαμεν παρὰ τοῦ Θεοῦ ἡμῶν περὶ τούτου, καὶ ἐπήκουσεν ἡμῖν.
\par }{\PP \VS{24}Καὶ διέστειλα ἀπὸ ἀρχόντων τῶν ἱερέων δώδεκα, τῷ Σαραΐα, τῷ Ἀσαβία, καὶ μετʼ αὐτῶν ἀπὸ ἀδελφῶν αὐτῶν δέκα.
\VS{25}Καὶ ἔστησα αὐτοῖς τὸ ἀργύριον καὶ τὸ χρυσίον καὶ τὰ σκεύη ἀπαρχῆς οἴκου Θεοῦ ἡμῶν, ἃ ὕψωσεν ὁ βασιλεὺς καὶ οἱ σύμβουλοι αὐτοῦ καὶ οἱ ἄρχοντες αὐτοῦ, καὶ πᾶς Ἰσραὴλ οἱ εὑρισκόμενοι.
\VS{26}Καὶ ἔστησα ἐπὶ χεῖρας αὐτῶν ἀργυρίου τάλαντα ἑξακόσια πεντήκοντα, καὶ σκεύη ἀργυρᾶ ἑκατὸν, καὶ τάλαντα χρυσίου ἑκατὸν,
\VS{27}καὶ χαφουρῆ χρυσοῖ εἴκοσι εἰς τὴν ὁδὸν χίλιοι, καὶ σκεύη χαλκοῦ στίλβοντος ἀγαθοῦ διάφορα ἐπιθυμητὰ ἐν χρυσίῳ.
\VS{28}Καὶ εἶπα πρὸς αὐτοὺς, ὑμεῖς ἅγιοι τῷ Κυρίῳ, καὶ τὰ σκεύη ἅγια, καὶ τὸ ἀργύριον καὶ τὸ χρυσίον ἑκούσια τῷ Κυρίῳ Θεῷ πατέρων ἡμῶν.
\VS{29}Ἀγρυπνεῖτε καὶ τηρεῖτε ἕως στῆτε ἐνώπιον ἀρχόντων τῶν ἱερέων καὶ τῶν Λευιτῶν καὶ τῶν ἀρχόντων τῶν πατριῶν ἐν Ἱερουσαλὴμ εἰς σκηνὰς οἴκου Κυρίου.
\VS{30}Καὶ ἐδέξαντο οἱ ἱερεῖς καὶ οἱ Λευῖται σταθμὸν τοῦ ἀργυρίου καὶ τοῦ χρυσίου καὶ τῶν σκευῶν, ἐνεγκεῖν εἰς Ἱερουσαλὴμ εἰς οἶκον Θεοῦ ἡμῶν.
\par }{\PP \VS{31}Καὶ ἐξῄραμεν ἀπὸ τοῦ ποταμοῦ τοῦ Ἀουὲ ἐν τῇ δωδεκάτῃ τοῦ μηνὸς τοῦ πρώτου τοῦ ἐλθεῖν εἰς Ἱερουσαλήμ· καὶ χεὶρ Θεοῦ ἡμῶν ἦν ἐφʼ ἡμῖν, καὶ ἐῤῥύσατο ἡμᾶς ἀπὸ χειρὸς ἐχθροῦ καὶ πολεμίου ἐν τῇ ὁδῷ.
\VS{32}Καὶ ἤλθομεν εἰς Ἱερουσαλὴμ, καὶ ἐκαθίσαμεν ἐκεῖ ἡμέρας τρεῖς.
\VS{33}Καὶ ἐγενήθη τῇ ἡμέρᾳ τῇ τετάρτῃ ἐστήσαμεν τὸ ἀργύριον καὶ τὸ χρυσίον καὶ τὰ σκεύη ἐν οἴκῳ Θεοῦ ἡμῶν ἐπὶ χεῖρα Μεριμὼθ υἱοῦ Οὐρία τοῦ ἱερέως, καὶ μετʼ αὐτοῦ Ἐλεάζαρ υἱὸς Φινεὲς, καὶ μετʼ αὐτῶν Ἰωζαβὰδ υἱὸς Ἰησοῦ, καὶ Νωαδία υἱὸς Βαναΐα οἱ Λευῖται.
\VS{34}Ἐν ἀριθμῷ καὶ ἐν σταθμῷ τὰ πάντα, καὶ ἐγράφη πᾶς ὁ σταθμός.
\par }{\PP \VS{35}Ἐν τῷ καιρῷ ἐκείνῳ οἱ ἐλθόντες ἐκ τῆς αἰχμαλωσίας υἱοὶ τῆς παροικίας, προσήνεγκαν ὁλοκαυτώσεις τῷ Θεῷ Ἰσραὴλ, μόσχους δώδεκα περὶ παντὸς Ἰσραὴλ, κριοὺς ἐννενηκονταὲξ, ἀμνοὺς ἑβδομηκονταεπτὰ, χιμάρους περὶ ἁμαρτίας δώδεκα, τὰ πάντα ὁλοκαυτώματα τῷ Κυρίῳ·
\VS{36}Καὶ ἔδωκαν τὸ νόμισμα τοῦ βασιλέως τοῖς διοικηταῖς τοῦ βασιλέως καὶ ἐπάρχοις πέραν τοῦ ποταμοῦ· καὶ ἐδόξασαν τὸν λαὸν καὶ τὸν οἶκον τοῦ Θεοῦ.

\par }\Chap{9}{\PP \VerseOne{1}Καὶ ὡς ἐτελέσθη ταῦτα, ἤγγισαν πρὸς μὲ οἱ ἄρχοντες, λέγοντες, οὐκ ἐχωρίσθη ὁ λαὸς Ἰσραὴλ καὶ οἱ ἱερεῖς καὶ οἱ Λευῖται ἀπὸ λαῶν τῶν γαιῶν ἐν μακρύμμασιν αὐτῶν, τῷ Χανανὶ ὁ Ἐθὶ, ὁ Φερεζὶ, ὁ Ἰεβουσὶ, ὁ Ἀμμωνὶ, ὁ Μωαβὶ, καὶ ὁ Μοσερὶ, καὶ ὁ Ἀμοῤῥὶ,
\VS{2}ὅτι ἐλάβοσαν ἀπὸ θυγατέρων αὐτῶν ἑαυτοῖς καὶ τοῖς υἱοῖς αὐτῶν· καὶ παρήχθη σπέρμα τὸ ἅγιον ἐν λαοῖς τῶν γαιῶν, καὶ χεὶρ τῶν ἀρχόντων ἐν τῇ ἀσυνθεσίᾳ ταύτῃ ἐν ἀρχῇ.
\VS{3}Καὶ ὡς ἤκουσα τὸν λόγον τοῦτον, διέῤῥηξα τὰ ἱμάτιά μου, καὶ ἐπαλλόμην, καὶ ἔτιλλον ἀπὸ τῶν τριχῶν τῆς κεφαλῆς μου καὶ ἀπὸ τοῦ πώγωνός μου, καὶ ἐκαθήμην ἠρεμάζων.
\VS{4}Καὶ συνήχθησαν πρὸς μὲ πᾶς ὁ διώκων λόγον Θεοῦ Ἰσραὴλ ἐπὶ ἀσυνθεσίᾳ τῆς ἀποικίας· κᾀγὼ καθήμενος ἠρεμάζων ἕως τῆς θυσίας τῆς ἑσπερινῆς.
\par }{\PP \VS{5}Καὶ ἐν θυσίᾳ τῇ ἑσπερινῇ ἀνέστην ἀπὸ ταπεινώσεώς μου· καὶ ἐν τῷ διαῤῥήξαί με τὰ ἱμάτιά μου, καὶ ἐπαλλόμην, καὶ κλίνω ἐπὶ τὰ γόνατά μου, καὶ ἐκπετάζω τὰς χεῖράς μου πρὸς Κύριον τὸν Θεὸν,
\VS{6}καὶ εἶπα, Κύριε, ἠσχύνθην καὶ ἐνετράπην τοῦ ὑψῶσαι, Θεέ μου, τὸ πρόσωπόν μου πρὸς σὲ, ὅτι αἱ ἀνομίαι ἡμῶν ἐπληθύνθησαν ὑπὲρ κεφαλῆς ἡμῶν, καὶ αἱ πλημμέλειαι ἡμῶν ἐμεγαλύνθησαν ἕως εἰς τὸν οὐρανόν.
\VS{7}Ἀπὸ ἡμερῶν πατέρων ἡμῶν ἐσμὲν ἐν πλημμελείᾳ μεγάλῃ ἕως τῆς ἡμέρας ταύτης· καὶ ἐν ταῖς ἀνομίαις ἡμῶν παρεδόθημεν ἡμεῖς καὶ οἱ βασιλεῖς ἡμῶν καὶ οἱ υἱοὶ ἡμῶν ἐν χειρὶ βασιλέων τῶν ἐθνῶν ἐν ῥομφαίᾳ, καὶ ἐν αἰχμαλωσίᾳ, καὶ ἐν διαρπαγῇ, καὶ ἐν αἰσχύνῃ προσώπου ἡμῶν, ὡς ἡ ἡμέρα αὕτη.
\VS{8}Καὶ νῦν ἐπιεικεύσατο ἡμῖν ὁ Θεὸς ἡμῶν τοῦ καταλιπεῖν ἡμῖν εἰς σωτηρίαν, καὶ δοῦναι ἡμῖν στήριγμα ἐν τόπῳ ἁγιάσματος αὐτοῦ, τοῦ φωτίσαι ὀφθαλμοὺς ἡμῶν, καὶ δοῦναι ζωοποίησιν μικρὰν ἐν τῇ δουλείᾳ ἡμῶν.
\VS{9}Ὅτι δοῦλοι ἐσμὲν, καὶ ἐν τῇ δουλείᾳ ἡμῶν οὐκ ἐγκατέλιπεν ἡμᾶς Κύριος ὁ Θεὸς ἡμῶν· καὶ ἔκλινεν ἐφʼ ἡμᾶς ἔλεος ἐνώπιον βασιλέων Περσῶν, δοῦναι ἡμῖν ζωοποίησιν τοῦ ὑψῶσαι αὐτοὺς τὸν οἶκον τοῦ Θεοῦ ἡμῶν, καὶ ἀναστῆσαι τὰ ἔρημα αὐτῆς, καὶ τοῦ δοῦναι ἡμῖν φραγμὸν ἐν Ἰούδα καὶ Ἱερουσαλήμ.
\VS{10}Τί εἴπωμεν ὁ Θεὸς ἡμῶν μετὰ τοῦτο; ὅτι ἐγκατελίπομεν ἐντολάς σου,
\VS{11}ἃς ἔδωκας ἡμῖν ἐν χειρὶ δούλων σου τῶν προφητῶν, λέγων, ἡ γῆ εἰς ἣν εἰσπορεύεσθε κληρονομῆσαι αὐτὴν, γῆ μετακινουμένη ἐστὶν ἐν μετακινήσει λαῶν τῶν ἐθνῶν ἐν μακρύμμασιν αὐτῶν, ὧν ἔπλησαν αὐτὴν ἀπὸ στόματος ἐπὶ στόμα ἐν ἀκαθαρσίαις αὐτῶν.
\par }{\PP \VS{12}Καὶ νῦν τὰς θυγατέρας ὑμῶν μὴ δότε τοῖς υἱοῖς αὐτῶν, καὶ ἀπὸ τῶν θυγατέρων αὐτῶν μὴ λάβητε τοῖς υἱοῖς ὑμῶν, καὶ οὐκ ἐκζητήσετε εἰρήνην αὐτῶν καὶ ἀγαθὸν αὐτῶν ἕως αἰῶνος, ὅπως ἐνισχύσητε, καὶ φάγητε τὰ ἀγαθὰ τῆς γῆς, καὶ κληροδοτήσητε τοῖς υἱοῖς ὑμῶν ἕως αἰῶνος.
\VS{13}Καὶ μετὰ πᾶν τὸ ἐρχόμενον ἐφʼ ἡμᾶς ἐν ποιήμασιν ἡμῶν τοῖς πονηροῖς καὶ ἐν πλημμελείᾳ ἡμῶν τῇ μεγάλῃ, ὅτι οὐκ ἔστιν ὡς ὁ Θεὸς ἡμῶν, ὅτι ἐκούφισας ἡμῶν τὰς ἀνομίας, καὶ ἔδωκας ἡμῖν σωτηρίαν·
\VS{14}Ὅτι ἐπεστρέψαμεν διασκεδάσαι ἐντολάς σου, καὶ ἐπιγαμβρεῦσαι τοῖς λαοῖς τῶν γαιῶν· μὴ παροξυνθῇς ἐν ἡμῖν ἕως συντελείας, τοῦ μὴ εἶναι ἐγκατάλειμμα καὶ διασωζόμενον.
\VS{15}Κύριε ὁ Θεὸς Ἰσραὴλ, δίκαιος σὺ, ὅτι κατελείφθημεν διασωζόμενοι, ὡς ἡ ἡμέρα αὕτη· ἰδοὺ ἡμεις ἐναντίον σου ἐν πλημμελείαις ἡμῶν, ὅτι οὐκ ἔστι στῆναι ἐνώπιόν σου ἐπὶ τούτῳ.

\par }\Chap{10}{\PP \VerseOne{1}Καὶ ὡς προσηύξατο Ἔσδρας, καὶ ὡς ἐξηγόρευσε κλαίων καὶ προσευχόμενος ἐνώπιον οἴκου τοῦ Θεοῦ, συνήχθησαν πρὸς αὐτὸν ἀπὸ Ἰσραὴλ ἐκκλησία πολλὴ σφόδρα, ἄνδρες καὶ γυναῖκες καὶ νεανίσκοι, ὅτι ἔκλαυσαν ὁ λαὸς, καὶ ὕψωσε κλαίων.
\VS{2}Καὶ ἀπεκρίθη Σεχενίας υἱὸς Ἰεὴλ ἀπὸ υἱῶν Ἠλὰμ, καὶ εἶπε τῷ Ἔσδρᾳ, ἡμεῖς ἠσυνθετήσαμεν τῷ Θεῷ ἡμῶν, καὶ ἐκαθίσαμεν γυναῖκας ἀλλοτρίας ἀπὸ τῶν λαῶν τῆς γῆς· καὶ νῦν ἐστιν ὑπομονὴ τῷ Ἰσραὴλ ἐπὶ τούτῳ.
\VS{3}Καὶ νῦν διαθώμεθα διαθήκην τῷ Θεῷ ἡμῶν ἐκβαλεῖν πάσας τὰς γυναῖκας, καὶ τὰ γενόμενα ἐξ αὐτῶν, ὡς ἂν βούλῃ· ἀνάστηθι, καὶ φοβέρισον αὐτοὺς ἐν ἐντολαῖς Θεοῦ ἡμῶν, καὶ ὡς ὁ νόμος, γενηθήτω.
\VS{4}Ἀνάστα, ὅτι ἐπὶ σὲ τὸ ῥῆμα, καὶ ἡμεῖς μετὰ σοῦ· κραταιοῦ καὶ ποίησον.
\par }{\PP \VS{5}Καὶ ἀνέστη Ἔσδρας, καὶ ὥρκισε τοὺς ἄρχοντας, τοὺς ἱερεῖς, καὶ Λευείτας, καὶ πάντα Ἰσραὴλ, τοῦ ποιῆσαι κατὰ τὸ ῥῆμα τοῦτο· καὶ ὤμοσαν.
\VS{6}Καὶ ἀνέστη Ἔσδρας ἀπὸ προσώπου οἴκου τοῦ Θεοῦ, καὶ ἐπορεύθη εἰς γαζοφυλάκιον Ἰωανὰν υἱοῦ Ἐλισοὺβ, και ἐπορεύθη ἐκεῖ· ἄρτον οὐκ ἔφαγε, καὶ ὕδωρ οὐκ ἔπιεν, ὅτι ἐπένθει ἐπὶ τῇ ἀσυνθεσίᾳ τῆς ἀποικίας.
\VS{7}Καὶ παρήνεγκαν φωνὴν ἐν Ἰούδα καὶ ἐν Ἱερουσαλὴμ πᾶσι τοῖς υἱοῖς τῆς ἀποικίας, τοῦ συναθροισθῆναι εἰς Ἱερουσαλήμ.
\VS{8}Πᾶς ὃς ἂν μὴ ἔλθῃ εἰς τρεῖς ἡμέρας, ὡς ἡ βουλὴ τῶν ἀρχόντων καὶ τῶν πρεσβυτέρων, ἀναθεματισθήσεται πᾶσα ἡ ὕπαρξις αὐτοῦ, καὶ αὐτὸς διασταλήσεται ἀπὸ ἐκκλησίας τῆς ἀποικίας.
\par }{\PP \VS{9}Καὶ συνήχθησαν πάντες ἄνδρες Ἰούδα καὶ Βενιαμὶν εἰς Ἱερουσαλὴμ εἰς τὰς τρεῖς ἡμέρας· οὗτος ὁ μὴν ὁ ἔννατος· ἐν εἰκάδι τοῦ μηνὸς ἐκάθισε πᾶς ὁ λαὸς ἐν πλατείᾳ οἴκου τοῦ Θεοῦ ἀπὸ θορύβου αὐτῶν περὶ τοῦ ῥήματος, καὶ ἀπὸ τοῦ χειμῶνος.
\VS{10}Καὶ ἀνέστη Ἔσδρας ὁ ἱερεὺς, καὶ εἶπε πρὸς αὐτοὺς, ὑμεῖς ἠσυνθετήκατε, καὶ ἐκαθίσατε γυναῖκας ἀλλοτρίας τοῦ προσθεῖναι ἐπὶ πλημμέλειαν Ἰσραήλ.
\VS{11}Καὶ νῦν δότε αἴνεσιν Κυρίῳ Θεῷ τῶν πατέρῶν ἡμῶν, καὶ ποιήσατε τὸ ἀρεστὸν ἐνώπιον αὐτοῦ, καὶ διαστάλητε ἀπὸ λαῶν τῆς γῆς καὶ ἀπὸ τῶν γυναικῶν τῶν ἀλλοτρίων.
\par }{\PP \VS{12}Καὶ ἀπεκρίθησαν πᾶσα ἡ ἐκκλησία, καὶ εἶπαν, μέγα τοῦτο τὸ ῥῆμά σου ἐφʼ ἡμᾶς ποιῆσαι.
\VS{13}Ἀλλὰ ὁ λαὸς πολὺς, καὶ ὁ καιρὸς χειμερινὸς, καὶ οὐκ ἔστι δύναμις στῆναι ἔξω· καὶ τὸ ἔργον οὐκ εἰς ἡμέραν μίαν καὶ οὐκ εἰς δύο, ὅτι ἐπληθύναμεν τοῦ ἀδικῆσαι ἐν τῷ ῥῆματι τούτῳ·
\VS{14}Στήτωσαν δὴ ἄρχοντες ἡμῶν, καὶ πᾶσι τοῖς ἐν πόλεσιν ἡμῶν ὃς ἐκάθισε γυναῖκας ἀλλοτρίας, ἐλθέτωσαν εἰς καιροὺς ἀπὸ συνταγῶν, καὶ μετʼ αὐτῶν πρεσβύτεροι πόλεως καὶ πόλεως, καὶ κριταὶ, τοῦ ἀποστρέψαι ὀργὴν θυμοῦ Θεοῦ ἡμῶν ἐξ ἡμῶν, περὶ τοῦ ῥήματος τούτου.
\VS{15}Πλὴν Ἰωνάθαν υἱὸς Ἀσὴλ, καὶ Ἰαζίας υἱὸς Θεκωὲ μετʼ ἐμοῦ περὶ τούτου· καὶ Μεσουλλὰμ, καὶ Σαββαθαῒ ὁ Λευίτης βοηθῶν αὐτοῖς.
\par }{\PP \VS{16}Καὶ ἐποίησαν οὕτως υἱοὶ τῆς ἀποικίας· Καὶ διεστάλησαν Ἔσρας ὁ ἱερεὺς, καὶ ἄνδρες ἄρχοντες πατριῶν τῷ οἴκῳ, καὶ πάντες ἐν ὀνόμασιν, ὅτι ἐπέστρεψαν ἐν ἡμέρᾳ μιᾷ τοῦ μηνὸς τοῦ δεκάτου ἐκζητῆσαι τὸ ῥῆμα.
\VS{17}Καὶ ἐτέλεσαν ἐν πᾶσιν ἀνδράσιν οἳ ἐκάθισαν γυναῖκας ἀλλοτρίας, ἕως ἡμέρας μιᾶς τοῦ μηνὸς τοῦ πρώτου.
\par }{\PP \VS{18}Καὶ εὑρέθη ἀπὸ υἱῶν τῶν ἱερέων οἳ ἐκάθισαν γυναῖκας ἀλλοτρίας, ἀπὸ υἱῶν Ἰησοῦ υἱοῦ Ἰωσεδὲκ, καὶ ἀδελφοὶ αὐτοῦ Μαασία, καὶ Ἐλιέζερ, καὶ Ἰαρὶβ, καὶ Γαδαλία.
\VS{19}Καὶ ἔδωκαν χεῖρα αὐτῶν τοῦ ἐξενέγκαι γυναῖκας ἑαυτῶν, καὶ πλημμελείας κριὸν ἐκ προβάτων περὶ πλημμελήσεως αὐτῶν.
\VS{20}Καὶ ἀπὸ υἱῶν Ἐμμὴρ, Ἀνανὶ, καὶ Ζαβδία.
\VS{21}Καὶ ἀπὸ υἱῶν Ἠρὰμ, Μασαὴλ, καὶ Ἐλία, καὶ Σαμαΐα, καὶ Ἰεὴλ, καὶ Ὀζία.
\VS{22}Καὶ ἀπὸ υἱῶν Φασοὺρ, Ἐλιωναῒ, Μαασία, καὶ Ἰσμαὴλ, καὶ Ναθαναὴλ, καὶ Ἰωζαβὰδ, καὶ Ἠλασά.
\VS{23}Καὶ ἀπὸ τῶν Λευιτῶν, Ἰωζαβὰδ, καὶ Σαμοὺ, καὶ Κωλία, αὐτὸς Κωλίτας, καὶ Φεθεΐα, καὶ Ἰούδας, καὶ Ἐλιέζερ.
\VS{24}Καὶ ἀπὸ τῶν ᾀδόντων, Ἐλισάβ· καὶ ἀπὸ τῶν πυλωρῶν, Σολμὴν, καὶ Τελμὴν, καὶ Ὠδούθ.
\VS{25}Καὶ ἀπὸ Ἰσραὴλ, ἀπὸ υἱῶν Φόρος, Ῥαμία, καὶ Ἀζία, καὶ Μελχία, καὶ Μεαμὶν, καὶ Ἐλεάζαρ, καὶ Ἀσαβία, καὶ Βαναία.
\VS{26}Καὶ ἀπὸ υἱῶν Ἡλὰμ, Ματθανία, καὶ Ζαχαρία, καὶ Ἰαϊὴλ, καὶ Ἀβδία, καὶ Ἰαριμὼθ, καὶ Ἠλία.
\VS{27}Καὶ ἀπὸ υἱῶν Ζαθούα, Ἐλιωναῒ, Ἐλισοὺβ, Ματθαναῒ, καὶ Ἀρμὼθ, καὶ Ζαβὰδ, καὶ Ὀζιζά.
\VS{28}Καὶ ἀπὸ υἱῶν Βαβεῒ, Ἰωανὰν, Ἀνανία, καὶ Ζαβοὺ, καὶ Θαλί.
\VS{29}Καὶ ἀπὸ υἱῶν Βανουῒ, Μοσολλὰμ, Μαλοὺχ, Ἀδαΐας, Ἰασοὺβ, καὶ Σαλουΐα, καὶ Ῥημώθ.
\VS{30}Καὶ ἀπὸ υἱῶν Φαὰθ Μωὰβ, Ἐδνὲ, καὶ Χαλὴλ, καὶ Βαναία, Μαασία, Ματθανία, Βεσελεὴλ, καὶ Βανουῒ, καὶ Μανασσῆ·
\VS{31}Καὶ ἀπὸ υἱῶν Ἠρὰμ, Ἐλιέζερ, Ἰεσία, Μελχία, Σαμαΐας, Σεμεὼν,
\VS{32}Βενιαμὶν, Βαλοὺχ, Σαμαρία.
\VS{33}Καὶ ἀπὸ υἱῶν Ἀσὴμ, Μετθανία, Ματθαθὰ, Ζαδὰβ, Ἐλιφαλὲτ, Ἱεραμὶ, Μανασσῆ, Σεμεΐ.
\VS{34}Καὶ ἀπὸ υἱῶν Βανὶ, Μοοδία, Ἀμρὰμ, Οὐὴλ,
\VS{35}Βαναία, Βαδαία, Χελκία,
\VS{36}Οὐουανία, Μαριμὼθ, Ἐλιασὶφ,
\VS{37}Ματθανία, Ματθαναΐ· καὶ ἐποίησαν
\VS{38}οἱ υἱοἱ Βανουὶ, καὶ οἱ υἱοὶ Σεμεῒ,
\VS{39}καὶ Σελεμία, καὶ Νάθαν, καὶ Ἀδαΐα,
\VS{40}Μαχαδναβοὺ, Σεσεῒ, Σαριοὺ,
\VS{41}Ἐζριὴλ, καὶ Σελεμία, καὶ Σαμαρία,
\VS{42}καὶ Σελλοὺμ, Ἀμαρεία, Ἰωσήφ.
\VS{43}Ἀπὸ υἱῶν Ναβοὺ, Ἰαὴλ, Ματθανίας, Ζαβὰδ, Ζεβεννὰς, Ἰαδαὶ, καὶ Ἰωὴλ, καὶ Βαναία.
\par }{\PP \VS{44}Πάντες οὗτοι ἐλάβοσαν γυναῖκας ἀλλοτρίας, καὶ ἐγέννησαν ἐξ αὐτῶν υἱούς.
\par }